% --- Section: Conclusion ---
% Drafted Phase 6 Checkpoint C. Merged with prior §07 threats-to-validity
% and future-work content per A.5 structure migration.
% Target: ~400 words (200 conclusion + 200 threats/limitations).

\section{Discussion and Conclusion}
\label{sec:conclusion}

We benchmarked five frontier LLMs on 171 Bayesian and inferential
statistical tasks with 473 perturbations, scored by a deterministic
keyword rubric, an NMACR aggregator with literature-derived weights,
and an out-of-family Llama 3.3 70B Turbo judge constrained to
assumption-compliance. The five models produce three different rank
orders depending on the axis. The model with the lowest base
accuracy (GPT-4.1, $0.621$) is the most robust to perturbation
($0.9996$); the model that leads accuracy ($0.683$) and binary-pass
calibration (ECE $0.0334$) drops to fourth on robustness ($0.9565$).
Independently, $44.3\%$ of responses do not articulate the
assumptions the derivation requires, and the keyword rubric
overcounts assumption articulation by $15.4$ percentage points
relative to the judge. A single number is the wrong instrument for
statistical-reasoning evaluation; the right ranking depends on the
property being measured.

\textbf{Threats to validity.} The Llama 3.3 70B Turbo judge can
carry systematic biases from its training. We constrain it to one
dimension rather than asking for a holistic preference signal, and
we validate it against the keyword rubric (Spearman $\rho = 0.602$
on assumption-compliance). Krippendorff's $\alpha$ on the three
judged dimensions sits in the ``questionable'' range. That reflects
genuine ambiguity in scoring free-form Bayesian reasoning, not
instrument failure, and we report it as honest disclosure. The
three-ranking divergence is a within-model cross-metric finding
and does not depend on inter-rater agreement. Task coverage is
broad for Bayesian and inferential reasoning but is not
representative of statistical reasoning generally; causal inference
and experimental design are out of scope. Mean perturbation count
is $2.77$ per base task, and ECE estimates use 10{,}000-sample
bootstrap CIs that carry the uncertainty through. Model versions
are May 2026 snapshots; the inversion patterns may not replicate
as APIs evolve.

\textbf{Future work.} Three extensions follow. First, broaden the
task suite into causal inference, experimental design, and
frequentist hypothesis testing, to see whether the three-axis
divergence generalizes outside Bayesian reasoning. Second,
replicate with a second out-of-family judge to bound judge-bias
variance; cross-judge agreement on assumption-compliance would
also act as an upper bound on inter-rater reliability for the
dimension. Third, treat the GPT-4.1 inversion as a hypothesis
about training-data composition and decoding choices, and run
controlled experiments at the provider level that isolate the
relevant factors.
